\documentclass[10pt]{article}
\usepackage[a4paper,margin=0.72in]{geometry}
\usepackage{amsmath,amssymb,amsthm}
\usepackage{xcolor}
\usepackage[hidelinks]{hyperref}

\newtheorem{theorem}{Theorem}
\newtheorem{lemma}[theorem]{Lemma}
\newtheorem{proposition}[theorem]{Proposition}
\newcommand{\F}{\mathbb F}
\newcommand{\REL}{\textup{REL}}

\title{Half-sparse irreducible polynomials over $\F_2$}
\author{Axeyum Lemire proof project}
\date{Working proof manuscript --- 20 August 2026}

\begin{document}
\maketitle
\vspace{-1.2em}
\begin{center}
\fcolorbox{red}{red!6}{\parbox{0.93\linewidth}{\centering
\textbf{Not yet a proof.}  Lemma~\ref{lem:top} is the sole open mathematical
statement.  Every other implication below is proved and replayed in Axeyum.
This notice must be removed only after Lemma~\ref{lem:top} has a checked proof.}}
\end{center}

\begin{abstract}
Kaser and Lemire asked whether, for every $n$, there is an irreducible
$f\in\F_2[x]$ of degree $n$ for which
$\deg(f-x^n)\leq\lfloor n/2\rfloor$.  We reduce the question to a single
one one-sided relative trace estimate for binary von Mangoldt populations.
The trace retains cancellation across only $\lceil\log_2\ell\rceil+2$ top
conductor levels and asks for a linear saving over the relative
Riemann--Hypothesis bound.  An independently checked finite calculation
handles $n\leq400$.
\end{abstract}

\section{The ray class and its binary refinements}

Put $\ell=\lceil n/2\rceil-1$ and
\[
 E_j=(\F_2[x]/x^{j+1})^*,\qquad |E_j|=2^j.
\]
For a monic polynomial $F$, write
$\langle F\rangle_j=x^{\deg F}F(x^{-1})\pmod{x^{j+1}}$ and define
\[
 N_j(b)=\sum_{\substack{F\ \mathrm{monic},\ \deg F=n\\
                         \langle F\rangle_j=b}}\Lambda(F)
 \quad(b\in E_j).
\]
The desired shape is equivalent to $\langle f\rangle_\ell=1$.  Thus it is
enough to prove that the identity class contains more von Mangoldt mass than
can be supplied by proper prime powers.

The quotient $E_j\to E_{j-1}$ has two-element fibres.  For a parent
$b\in E_{j-1}$, let $b_0=b$ be its degree-$<j$ lift and
$b_1=b_0(1+x^j)$, and put
\[
 H_j(b)=N_j(b_0)-N_j(b_1),\qquad H_j^*=\max_b|H_j(b)|.
\]

\begin{proposition}[Haar reconstruction]\label{prop:haar}
For every $e\in E_\ell$,
\[
 2^\ell N_\ell(e)
 =2^n+\sum_{j=1}^{\ell}\varepsilon_j(e)2^{j-1}H_j(b_j(e)),
 \qquad \varepsilon_j(e)\in\{-1,1\}.
\]
Consequently
\begin{equation}\label{eq:triangle}
 \sum_{j=1}^{\ell}2^{j-1}H_j^*<2^{2\ell}
 \quad\Longrightarrow\quad
 |N_\ell(e)-2^{n-\ell}|<2^\ell\quad(e\in E_\ell).
\end{equation}
\end{proposition}

\begin{proof}
At each binary split, a child's value is half the parent value plus or minus
half the sibling difference.  Iterating from
$N_0(1)=\sum_{\deg F=n}\Lambda(F)=2^n$ gives the identity.  Taking absolute
values gives \eqref{eq:triangle}.
\end{proof}

Let $X_j$ be the $2^{j-1}$ characters of $E_j$ that do not factor through
$E_{j-1}$, and set
$S_n(\chi)=\sum_{\deg F=n}\Lambda(F)\chi(\langle F\rangle_j)$.
Fourier inversion on a two-child fibre gives
\begin{equation}\label{eq:fourier}
 H_j(b)=2^{1-j}\sum_{\chi\in X_j}\overline{\chi(b)}S_n(\chi).
\end{equation}
The $L$-polynomial of such a character has degree $j-1$.  The function-field
Riemann hypothesis and \eqref{eq:fourier} therefore give
\begin{equation}\label{eq:weil}
 H_j^*\leq (j-1)2^{\lceil n/2\rceil}.
\end{equation}

\section{The one remaining estimate}

Put $c=\lceil\log_2\ell\rceil$, $a=\ell-c-1$, and
\[
 W_{\ell,n}=2^{\lceil n/2\rceil}
       \sum_{1\leq j<a}(j-1)2^{j-1},\qquad
 B_{\ell,n}=2^{2\ell}-W_{\ell,n}.
\]
For $\ell\geq200$ the elementary geometric sum gives $B_{\ell,n}>0$.
Along the identity path the top increments telescope:
\begin{equation}\label{eq:connected}
 C_{\ell,n}:=\sum_{j=a}^{\ell}2^{j-1}H_j(1)
 =2^\ell N_\ell(1)-2^{a-1}N_{a-1}(1).
\end{equation}

\begin{lemma}[One-sided relative trace]\label{lem:top}
For $\ell\geq200$ and $n\in\{2\ell+1,2\ell+2\}$,
\begin{equation}\tag{$\REL$}\label{eq:rel}
 C_{\ell,n}>-B_{\ell,n}.
\end{equation}
\end{lemma}

\noindent\textcolor{red}{\textbf{Proof status: open.}}
Equation~\eqref{eq:connected} is also the difference of the
$\F_{2^n}$-point counts on two levels of the Carlitz cyclotomic
tower.  The lemma is one-sided---an arbitrarily large positive trace is
allowed---and preserves cancellation across all $c+2$ retained levels.
At $\ell=200$ its separate-level Riemann--Hypothesis envelope exceeds the
allowed negative magnitude by a factor $626$; asymptotically the requested
saving is $4\ell+O(\log\ell)$.  This saving is needed only in character
families whose exact order exceeds the largest power $Q$ with
$3Q\lceil\log_2\ell\rceil\leq\ell$: the individual Riemann--Hypothesis bounds
for all smaller orders fit in the budget.  Thus only $O(\log\log\ell)$ order
bands per conductor remain, all genuinely high-Witt rather than ordinary
Artin--Schreier families.

\begin{proposition}[The closing calculation]\label{prop:close}
Lemma~\ref{lem:top} implies
\[
 N_\ell(1)>2^{n-\ell}-2^\ell
\]
for both endpoint degrees and every $\ell\geq200$.
\end{proposition}

\begin{proof}
The part of Proposition~\ref{prop:haar}'s identity below $a$ is at least
$-W_{\ell,n}$ by \eqref{eq:weil}.  Keeping the remaining levels together and
using \eqref{eq:rel} therefore gives
\[
 2^\ell N_\ell(1)-2^n
 >-W_{\ell,n}-B_{\ell,n}=-2^{2\ell}.
\]
Division by $2^\ell$ proves the assertion.
\end{proof}

\section{Completion of the conjecture}

\begin{lemma}[Proper powers]\label{lem:powers}
For $n=2\ell+1$, one has
$N_\ell(1)=1+nI_n(1)$, where $I_n(1)$ counts degree-$n$ irreducibles in the
identity class.  For $n=2\ell+2$ and $\ell\geq200$, the total identity-class
mass of proper prime powers is strictly smaller than
$2^{n-\ell}-2^\ell$.
\end{lemma}

\begin{proof}
At the odd endpoint every proper divisor of $n$ is at most $\ell$; the only
irreducible of that degree in the identity ray class is $x$, and odd powering
is an automorphism of the $2$-group $E_\ell$.  This gives the first identity.
At the even endpoint, every odd exponent $k\geq3$ contributes zero by the
same automorphism argument.  Squares reduce to the $2$-torsion subgroup,
of order $2^{\lceil\ell/2\rceil}$; all remaining even exponents are at least
four.  Counting all monic irreducibles at their possible degrees (and
including the ramified power $x^n$) gives the explicit upper bound
\[
 (\ell+1)2^{\lceil\ell/2\rceil}
 +(2\ell+2)2^{\lceil(\ell+1)/2\rceil},
\]
which is smaller than $2^{\ell+2}-2^\ell$ for $\ell\geq200$.
\end{proof}

\begin{theorem}\label{thm:lemire}
Assuming Lemma~\ref{lem:top}, for every $n\geq1$ there is a monic irreducible
$f\in\F_2[x]$ of degree $n$ satisfying
$\deg(f-x^n)\leq\lfloor n/2\rfloor$.
\end{theorem}

\begin{proof}
Degrees $n\leq400$ are covered by independently checked Rabin certificates.
For $n>400$, Proposition~\ref{prop:close} and
Proposition~\ref{prop:haar} give
$N_\ell(1)>2^{n-\ell}-2^\ell$.  Lemma~\ref{lem:powers} shows that this exceeds
the proper-prime-power mass, so an irreducible lies in the identity class.
Taking reciprocals gives the asserted half-sparse polynomial.
\end{proof}

\begin{thebibliography}{9}\small
\bibitem{KL} O.~Kaser and D.~Lemire,
\emph{Strongly universal string hashing is fast},
The Computer Journal \textbf{57} (2014), 1624--1638.
\bibitem{Hayes} D.~R.~Hayes,
\emph{The distribution of irreducibles in $\mathrm{GF}[q,x]$},
Trans. Amer. Math. Soc. \textbf{117} (1965), 101--127.
\end{thebibliography}

\end{document}
